\documentclass[11pt]{article}

\usepackage[margin=1in]{geometry}
\usepackage[T1]{fontenc}
\usepackage{lmodern}
\usepackage{microtype}
\usepackage{amsmath,amssymb,mathtools}
\usepackage{booktabs}
\usepackage{array}
\usepackage{enumitem}
\usepackage{xcolor}
\usepackage{hyperref}
\usepackage[nameinlink,noabbrev]{cleveref}

\hypersetup{
  colorlinks=true,
  linkcolor=blue!55!black,
  citecolor=blue!55!black,
  urlcolor=blue!55!black,
  pdftitle={Background for Algorithm 5.1 of Elsenhans--Jahnel},
  pdfauthor={OpenAI}
}

\setlist[enumerate]{itemsep=0.35em,topsep=0.45em}
\setlist[itemize]{itemsep=0.25em,topsep=0.4em}
\setlength{\parindent}{0pt}
\setlength{\parskip}{0.55em}

\newcommand{\QQ}{\mathbf{Q}}
\newcommand{\ZZ}{\mathbf{Z}}
\newcommand{\PP}{\mathbf{P}}
\newcommand{\AA}{\mathbf{A}}
\newcommand{\GG}{\mathbf{G}}
\newcommand{\ol}[1]{\overline{#1}}
\newcommand{\Gal}{\operatorname{Gal}}
\newcommand{\Tr}{\operatorname{Tr}}
\newcommand{\Pic}{\operatorname{Pic}}
\newcommand{\Spec}{\operatorname{Spec}}
\newcommand{\Sym}{\operatorname{Sym}}
\newcommand{\Span}{\operatorname{Span}}
\newcommand{\Mat}{\operatorname{Mat}}
\newcommand{\im}{\operatorname{im}}
\newcommand{\rank}{\operatorname{rank}}
\newcommand{\disc}{\operatorname{disc}}
\newcommand{\cM}{\mathcal{M}}
\newcommand{\tcM}{\widetilde{\mathcal{M}}}
\newcommand{\tM}{\widetilde{M}}
\newcommand{\cR}{\mathcal{R}}
\newcommand{\cA}{\mathscr{A}}
\newcommand{\cL}{\mathcal{L}}

\newenvironment{roadmap}{%
  \begin{center}\begin{minipage}{0.93\textwidth}\small
  \hrule\vspace{0.55em}\textbf{Logical role.}\ }{%
  \vspace{0.55em}\hrule\end{minipage}\end{center}}

\title{Background Needed to Understand and Implement\\
Algorithm 5.1 of Elsenhans--Jahnel}
\author{A detailed, proof-free guide to arXiv:1209.5591v3}
\date{\today}

\begin{document}
\maketitle

\begin{abstract}
Algorithm 5.1 in Elsenhans--Jahnel constructs a smooth cubic surface over
$\QQ$ whose Galois action on the $27$ lines realizes a prescribed subgroup
of the Weyl group $W(E_6)$.  The algorithm combines four ingredients:
Coble's projective coordinates on the moduli space of marked cubic surfaces,
a linear action of $W(E_6)$ on those coordinates, explicit Galois descent of
a twisted moduli space, and a second descent that reconstructs an equation
from Clebsch's invariants.  This note collects the definitions and the
sequence of constructions needed to understand those ingredients.  It does
not reproduce the proofs in the paper.  The emphasis is on the objects that
must actually be represented in an implementation and on the maps that
connect one stage to the next.
\end{abstract}

\tableofcontents

\section{What the algorithm does}

Fix a subgroup
\[
  G\subseteq W(E_6)
\]
and a finite Galois extension $L/\QQ$ together with an explicit isomorphism
\[
  \rho:\Gal(L/\QQ)\xrightarrow{\sim}G.
\]
Algorithm 5.1 seeks a smooth cubic surface $C/\QQ$ for which the action of
$\Gal(\ol\QQ/\QQ)$ on the $27$ geometric lines has image conjugate to the
specified subgroup $G\subseteq W(E_6)$.  The embedding of $G$ into $W(E_6)$,
not merely the abstract isomorphism type of $G$, is part of the input.

At a high level, the data pass through the following chain:
\[
\begin{gathered}
 (G,L,\rho)
 \longrightarrow
 \text{a twisted form }\tcM_{\rho}\text{ of the marked moduli space}
 \longrightarrow
 \tcM_{\rho}(\QQ) \\
 \longrightarrow
 \text{Coble coordinates }(\gamma_1,\ldots,\gamma_{40})
 \longrightarrow
 [A:B:C:D:E] \\
 \longrightarrow
 (\sigma_1,\ldots,\sigma_5)
 \longrightarrow
 g(T)=\prod_{i=0}^{4}(T-a_i)
 \longrightarrow
 \text{an explicit cubic equation over }\QQ.
\end{gathered}
\]

There are two conceptually distinct Galois descents.

\begin{enumerate}[label=\textbf{Descent \arabic*.}]
\item The first descent constructs the twist $\tcM_{\rho}$ and an explicit
projective model of it over $\QQ$.  This descent encodes the prescribed
Galois action on the $27$ lines.
\item The second descent starts from the five pentahedral coefficients
$a_0,\ldots,a_4$, usually defined only over a splitting field, and produces a
cubic equation over $\QQ$ by a trace construction.
\end{enumerate}

The input field $L$ in the first descent and the degree-five etale algebra
used in the second descent are different objects and need not be related.
This distinction prevents a common source of confusion.

\begin{roadmap}
The marked moduli space remembers enough information to prescribe the Galois
permutation of the lines.  Coble's coordinates make that moduli space
explicit.  Clebsch's invariants forget the marking.  The equation problem
then turns the unmarked invariants into a concrete cubic polynomial.
\end{roadmap}

\section{Cubic surfaces, markings, and the group $W(E_6)$}

\subsection{Smooth cubic surfaces and their lines}

A smooth cubic surface over a field $K$ is a smooth degree-three hypersurface
$C\subset \PP^3_K$.  Over an algebraic closure $\ol K$, every smooth cubic
surface contains exactly $27$ lines.  Incidence among these lines is
preserved by every automorphism of the surface and by the action of
$\Gal(\ol K/K)$.

The group of permutations of the $27$ lines preserving their intersection
pairing is naturally isomorphic to the Weyl group $W(E_6)$, of order
\[
  |W(E_6)|=51{,}840.
\]
Consequently, after choosing labels for the lines, a cubic surface over $K$
gives a homomorphism
\[
  \Gal(\ol K/K)\longrightarrow W(E_6).
\]
Changing the labeling conjugates the image.  Thus an unmarked cubic surface
determines a conjugacy class of subgroups of $W(E_6)$.

\subsection{The blow-up model and a marking}

A marking is an ordered six-tuple
\[
  (\ell_1,\ldots,\ell_6)
\]
of mutually disjoint lines on $C$.  Blowing down these six lines identifies
$C_{\ol K}$ with the blow-up of $\PP^2_{\ol K}$ at six ordered points
$p_1,\ldots,p_6$ in general position.  Here ``general position'' means:

\begin{itemize}
\item no three of the six points are collinear;
\item the six points do not all lie on a conic.
\end{itemize}

Write $H$ for the pullback of a line in $\PP^2$ and $E_i$ for the exceptional
curve over $p_i$.  The marking determines the standard labeling of all
$27$ lines:
\[
\begin{array}{lll}
 \ell_i=E_i, & 1\leq i\leq 6, & \text{six exceptional curves},\\[2mm]
 \ell_{ij}''=H-E_i-E_j, & 1\leq i<j\leq 6, & \text{fifteen transforms of lines},\\[2mm]
 \ell_i'=2H-\displaystyle\sum_{j\neq i}E_j,
 & 1\leq i\leq 6, & \text{six transforms of conics}.
\end{array}
\]
The orthogonal complement of the canonical class in $\Pic(C_{\ol K})$ is a
root lattice of type $E_6$.  The corresponding Weyl group acts by changing
the marking.

A useful generating set for $W(E_6)$ consists of:

\begin{itemize}
\item the subgroup $S_6$, which permutes the six blown-up points;
\item one quadratic Cremona transformation centered at, say,
$p_1,p_2,p_3$.
\end{itemize}

This description is important when one constructs the permutation action on
Coble's invariants explicitly.

\section{The two moduli spaces}

\subsection{The marked moduli space}

Let $U\subset (\PP^2)^6$ be the open set of ordered six-tuples of points in
general position.  The projective linear group $\operatorname{PGL}_3$ acts by
simultaneous change of coordinates.  The quotient
\[
  \tcM = U/\operatorname{PGL}_3
\]
is the fine moduli space of marked smooth cubic surfaces.  It is a smooth
affine fourfold.  The marking rigidifies the surface: an automorphism that
fixes the marked lines is the identity, so a fine moduli space can exist.

A convenient affine chart is obtained by sending the first four points to
\[
 (1:0:0),\quad (0:1:0),\quad (0:0:1),\quad (1:1:1).
\]
A marked surface can then be represented by the columns of
\[
 \begin{pmatrix}
  1&0&0&1&w&y\\
  0&1&0&1&x&z\\
  0&0&1&1&1&1
 \end{pmatrix},
\]
subject to explicit nonvanishing conditions expressing general position.
This chart is useful for evaluation and interpolation computations.

The group $W(E_6)$ acts on $\tcM$ by changing the marking.  Denote the
corresponding automorphism by
\[
  T_g:\tcM\longrightarrow\tcM,
  \qquad g\in W(E_6).
\]

\subsection{The unmarked moduli space}

The coarse moduli space of smooth cubic surfaces is
\[
  \cM=\tcM/W(E_6).
\]
The forgetful morphism
\[
  \operatorname{pr}:\tcM\longrightarrow\cM
\]
is generically a covering of degree $51{,}840$.  The unmarked moduli space is
not fine because some cubic surfaces have nontrivial automorphisms.

Clebsch's invariant theory gives an open embedding
\[
  \operatorname{Cl}:\cM\hookrightarrow \PP(1,2,3,4,5).
\]
The weighted homogeneous coordinates are denoted
\[
  [A:B:C:D:E],
\]
with weights $1,2,3,4,5$, respectively.

The algorithm uses the marked moduli space to encode the Galois action, but
it uses the unmarked invariants $[A:B:C:D:E]$ to recover an equation.

\section{Coble's irrational invariants}

\subsection{The basic determinants}

Represent the ordered points $p_i$ by row vectors
\[
  p_i=(x_{i,0}:x_{i,1}:x_{i,2}),\qquad 1\leq i\leq 6.
\]
For three distinct indices define the $3\times 3$ minor
\[
  m_{\{i,j,k\}}
  =\det\begin{pmatrix}
    x_{i,0}&x_{i,1}&x_{i,2}\\
    x_{j,0}&x_{j,1}&x_{j,2}\\
    x_{k,0}&x_{k,1}&x_{k,2}
  \end{pmatrix}.
\]
It vanishes exactly when $p_i,p_j,p_k$ are collinear.

Let $d_2$ be the determinant of the $6\times6$ matrix obtained by evaluating
the quadratic monomials at the six points:
\[
 d_2=
 \det\begin{pmatrix}
 x_{1,0}^2&x_{1,1}^2&x_{1,2}^2&x_{1,0}x_{1,1}&x_{1,0}x_{1,2}&x_{1,1}x_{1,2}\\
 \vdots&\vdots&\vdots&\vdots&\vdots&\vdots\\
 x_{6,0}^2&x_{6,1}^2&x_{6,2}^2&x_{6,0}x_{6,1}&x_{6,0}x_{6,2}&x_{6,1}x_{6,2}
 \end{pmatrix}.
\]
It vanishes when the six points lie on a conic.

\subsection{The forty sections}

Coble defines forty multihomogeneous invariants.  There are two types.

For a partition of $\{1,\ldots,6\}$ into two unordered triples, set
\[
 \gamma_{(i_1i_2i_3)(i_4i_5i_6)}
 =m_{\{i_1,i_2,i_3\}}m_{\{i_4,i_5,i_6\}}d_2.
\]
There are ten invariants of this type.

For a cyclically ordered triple of unordered pairs, set
\[
\begin{aligned}
 \gamma_{(i_1i_2)(i_3i_4)(i_5i_6)}={}&
 m_{\{i_1,i_3,i_4\}}m_{\{i_2,i_3,i_4\}}
 m_{\{i_3,i_5,i_6\}}m_{\{i_4,i_5,i_6\}}\\
 &\qquad\cdot
 m_{\{i_5,i_1,i_2\}}m_{\{i_6,i_1,i_2\}}.
\end{aligned}
\]
There are thirty invariants of this type.  A concrete implementation must
choose and retain a canonical ordering of these forty symbols.

Each $\gamma$ is of degree three in the coordinates of each point.  It is
invariant under the simultaneous $\operatorname{PGL}_3$ action.  On $U$, none
of the forty sections vanishes.  The adjective ``irrational'' is historical:
these are invariants of a \emph{marked} surface, not invariants under the full
$W(E_6)$ action.  Their values therefore usually live in a field over which
the marking is defined.

\subsection{The gamma map and its equations}

The forty sections define Coble's gamma map
\[
  \gamma:\tcM\hookrightarrow\PP^{39}.
\]
Although there are forty coordinate functions, their linear span has
dimension ten.  Equivalently, there are thirty independent linear relations
among them, and the image lies in a distinguished $\PP^9\subset\PP^{39}$.

Let $\tM$ denote the Zariski closure of the image.  This is Coble's gamma
variety.  Inside the $\PP^9$ determined by the linear relations, $\tM$ is the
common zero locus of a thirty-dimensional space of cubic equations.  Thus:
\[
 \dim \Sym^3(\QQ^{10})=220,
 \qquad
 \dim \cR_3=30,
\]
where $\cR_3$ denotes the space of cubic relations.  The variety has
dimension four and is not a complete intersection; the thirty cubics are
highly dependent as equations of a codimension-five variety.

One illustrative cubic relation is
\[
\begin{aligned}
&\gamma_{(12)(34)(56)}\gamma_{(23)(45)(16)}\gamma_{(14)(36)(25)}\\
&\qquad=
\gamma_{(12)(36)(45)}\gamma_{(34)(25)(16)}\gamma_{(56)(14)(23)}.
\end{aligned}
\]
Its $W(E_6)$-orbit spans the thirty-dimensional relation space.

For implementation, one chooses ten linearly independent signed gamma
sections
\[
  \eta_1,\ldots,\eta_{10}
\]
and expresses every signed invariant as a rational linear combination
\[
  \Gamma_a=\sum_{j=1}^{10}J_{aj}\eta_j,
  \qquad 1\leq a\leq80.
\]
Here $\Gamma_1,\ldots,\Gamma_{80}$ denotes the list
$\{\pm\gamma_1,\ldots,\pm\gamma_{40}\}$, and
\[
  J\in\Mat_{80\times10}(\QQ)
\]
is a fixed rank-ten matrix.  If $y=(y_1,\ldots,y_{10})^t$ is the vector of
values of the basis sections at a point, the full signed coordinate vector is
\[
  x=Jy\in L^{80}.
\]
The matrix $J$ and a basis of the thirty cubic relations are universal data:
they are computed once and then reused for every group and every number
field.

\section{The $W(E_6)$ action on Coble coordinates}

\subsection{Why eighty signed coordinates are used}

The line bundle defining the gamma map admits a $W(E_6)$-linearization for
which every $g\in W(E_6)$ sends each $\gamma_i$ to $\pm\gamma_j$.  Thus the
action on the forty coordinates is a signed permutation action.

To replace signed permutation matrices by ordinary permutations, the paper
uses the eighty sections
\[
  \pm\gamma_1,\ldots,\pm\gamma_{40}.
\]
They define an embedding
\[
  \gamma':\tcM\hookrightarrow\PP^{79}
\]
that is linearly equivalent to the original gamma map.  The image is still
contained in a $\PP^9$: the duplication introduces forty relations of the
form $(-\gamma_i)+\gamma_i=0$, in addition to the thirty relations among the
original forty gamma sections.

There is a faithful permutation representation
\[
  \Pi:W(E_6)\hookrightarrow S_{80}
\]
with forty blocks $\{\gamma_i,-\gamma_i\}$.

\subsection{Point-action convention and the inverse permutation}

Let $P_g$ be the $80\times80$ permutation matrix acting on a column of point
coordinates by
\[
  (P_gx)_a=x_{\Pi(g)^{-1}(a)}.
\]
Equivalently, the coordinate originally in position $a$ is moved to position
$\Pi(g)(a)$.  This explains the inverse permutation in the formula in
Remark 4.4 of the paper.

Since the ten-dimensional linear subspace is stable, there is a unique
matrix
\[
  R_g\in\operatorname{GL}_{10}(\QQ)
\]
satisfying
\[
  P_gJ=JR_g.
\]
The matrices $R_g$ are the ones needed for the actual descent computation.
They describe the action on the ten point-coordinate parameters $y$.

\begin{roadmap}
There are three equivalent representations of the same action:
\[
\begin{array}{ccl}
80\text{ coordinates} &:& \text{ordinary permutation matrices }P_g,\\
40\text{ coordinates} &:& \text{signed permutation matrices},\\
10\text{ coordinates} &:& \text{dense matrices }R_g.
\end{array}
\]
The paper uses the first model to state the descent cleanly.  An efficient
implementation normally performs the linear algebra in the third model, and
may use the second model for the universal precomputation.
\end{roadmap}

\section{Clebsch invariants and the map that forgets the marking}

Let
\[
  P_k=\sum_{i=1}^{40}\gamma_i^k.
\]
The sum is over the forty original gamma coordinates, one from each signed
pair.  The formulas below involve only even $k$, so the chosen sign of each
$\gamma_i$ is irrelevant.  One must not sum over all eighty signed
coordinates, since doing so inserts an unwanted factor of two in each
$P_k$.

Theorem 3.9 of the paper gives the following expressions for the Clebsch
invariants:
\[
\begin{aligned}
 A={}&-6P_2,\\
 B={}&-24P_4+\frac{41}{16}P_2^2,\\
 C={}&\frac{576}{13}P_6-\frac{396}{13}P_4P_2
       +\frac{29}{13}P_2^3,\\
 D={}&-\frac{62208}{1171}P_8
       +\frac{54864}{1171}P_6P_2
       +\frac{203616}{1171}P_4^2
       -\frac{61287}{1171}P_4P_2^2
       +\frac{13393}{4684}P_2^4,\\
 E={}&\frac{41472}{155}P_{10}
       -\frac{4605984}{36301}P_8P_2
       -\frac{106272}{403}P_6P_4
       +\frac{19990440}{471913}P_6P_2^2\\
   &\quad
       +\frac{47719206}{471913}P_4^2P_2
       -\frac{7468023}{471913}P_4P_2^3
       +\frac{10108327}{18876520}P_2^5.
\end{aligned}
\]

If all gamma coordinates are multiplied by $\lambda$, then the five
expressions scale by
\[
  (\lambda^2,\lambda^4,\lambda^6,\lambda^8,\lambda^{10}).
\]
This is exactly weighted projective scaling in $\PP(1,2,3,4,5)$, with
weighted scalar $\lambda^2$.  Hence the formulas are well defined on the
projective gamma point.

The expressions are $W(E_6)$-invariant.  Therefore, when the gamma
coordinates satisfy the twisted Galois compatibility condition, the resulting
Clebsch invariant vector is fixed by Galois and lies in
$\PP(1,2,3,4,5)(\QQ)$.

\section{Twisting the marked moduli space}

\subsection{The homomorphism that prescribes the line action}

Let $L/\QQ$ be finite Galois and let
\[
  \rho:\Gal(L/\QQ)\xrightarrow{\sim}G\subseteq W(E_6).
\]
Composing the restriction map
$\Gal(\ol\QQ/\QQ)\twoheadrightarrow\Gal(L/\QQ)$ with $\rho$ gives the desired
continuous homomorphism into $W(E_6)$.

Over $L$, the twisted moduli space becomes the ordinary marked moduli space:
\[
  (\tcM_{\rho})_L\cong\tcM_L.
\]
Its descent datum is
\[
  \delta_\sigma=T_{\rho(\sigma)}\circ\sigma,
  \qquad \sigma\in\Gal(L/\QQ),
\]
where $\sigma$ acts on coefficients and $T_{\rho(\sigma)}$ changes the
marking.

A $\QQ$-rational point of $\tcM_{\rho}$ is therefore represented by an
$L$-rational marked cubic surface whose Galois conjugation is compensated by
the prescribed relabeling.  After forgetting the marking, the underlying
cubic surface descends to $\QQ$, and the Galois action on its lines is the one
encoded by $\rho$.

\subsection{The projective fixed-point condition}

For a point
\[
  x=(x_1:\cdots:x_{80})\in\gamma'(\tcM_L)\subset\PP^{79}(L),
\]
the condition for it to define a $\QQ$-point on the twist is
\[
  P_{\rho(\sigma)}\sigma(x)=x
\]
in projective space.  In coordinates, this reads
\[
 \bigl(\sigma(x_{\Pi(\rho(\sigma))^{-1}(1)}):\cdots:
       \sigma(x_{\Pi(\rho(\sigma))^{-1}(80)})\bigr)
   =(x_1:\cdots:x_{80}).
\]

The paper then imposes exact equality on vectors rather than equality only
up to scalar:
\[
  P_{\rho(\sigma)}\sigma(x)=x\qquad\text{in }L^{80}.
\]
The fixed vectors form an $80$-dimensional $\QQ$-vector space in $L^{80}$.
Intersecting with the ten-dimensional linear subspace containing the gamma
variety gives a ten-dimensional $\QQ$-vector space $V$.  Galois descent gives
\[
  V\otimes_{\QQ}L\cong L^{10},
\]
and $\PP(V)$ is the descended ambient projective space.  Thus the stronger
vector condition is the mechanism for constructing the correct projective
form; it does not discard the projective descent.

\section{The explicit linear system defining $V$}

This section rewrites steps (i)--(iv) of Algorithm 5.1 in a matrix form suited
to implementation.

\subsection{Pairing group generators with field automorphisms}

Choose generators
\[
  \Gamma=\{g_1,\ldots,g_s\}\subseteq G.
\]
For each $g_i$, compute
\[
  \sigma_i=\rho^{-1}(g_i)\in\Gal(L/\QQ).
\]
The universal precomputation supplies $P_{g_i}$, $J$, and hence $R_{g_i}$ via
\[
  P_{g_i}J=JR_{g_i}.
\]

Choose a $\QQ$-basis
\[
  \omega_1,\ldots,\omega_n
  \qquad(n=[L:\QQ])
\]
of $L$.  Let $S_{\sigma_i}\in\operatorname{GL}_n(\QQ)$ be the matrix of
$\sigma_i$ in this basis, defined by
\[
  \sigma_i(\omega_b)=\sum_{a=1}^{n}(S_{\sigma_i})_{ab}\omega_a.
\]

\subsection{The ten-coordinate fixed equation}

Write $y=(y_1,\ldots,y_{10})^t\in L^{10}$.  Because $x=Jy$ and $J$ has rational
entries, the eighty-coordinate condition is equivalent to
\[
  R_{g_i}\,\sigma_i(y)=y,
  \qquad i=1,\ldots,s.
\]
This is the reduced form of the descent equations.

For a fully rational matrix formulation, write
\[
  y_j=\sum_{a=1}^{n}U_{aj}\omega_a
\]
with $U\in\Mat_{n\times10}(\QQ)$.  The equation above becomes
\[
  S_{\sigma_i}UR_{g_i}^{t}=U.
\]
After vectorization, this is
\[
  \bigl(R_{g_i}\otimes S_{\sigma_i}-I_{10n}\bigr)\operatorname{vec}(U)=0.
\]
Stack these systems for the chosen generators and compute their common
kernel.  Its $\QQ$-dimension should be ten.  A basis gives vectors
\[
  v_1,\ldots,v_{10}\in L^{10}.
\]
Put them as columns of a matrix
\[
  B_{\rho}=(v_1\ \cdots\ v_{10})\in\operatorname{GL}_{10}(L).
\]
Then
\[
  y=B_{\rho}z,
  \qquad z\in\QQ^{10},
\]
parametrizes the fixed vector space $V$.

The paper formulates the system in eighty coordinates and therefore counts
$80[L:\QQ]\#\Gamma$ equations in $10[L:\QQ]$ unknowns.  Passing first to the
ten-dimensional representation gives the same fixed space with many fewer
redundant equations.

\subsection{Why generators suffice}

The maps $y\mapsto R_{\rho(\sigma)}\sigma(y)$ form a semilinear group action.
Hence a vector fixed by the chosen generating automorphisms is fixed by the
entire Galois group.  No equations for all elements of $G$ are required.

\section{Descending the thirty cubic equations}

Choose a basis
\[
  F_1,\ldots,F_{30}\in\QQ[y_1,\ldots,y_{10}]_3
\]
of the universal cubic relation space of the gamma variety.  Substitute
\[
  y=B_{\rho}z.
\]
This produces cubic forms
\[
  F_a(B_{\rho}z)\in L[z_1,\ldots,z_{10}]_3.
\]
Expand each coefficient in the fixed $\QQ$-basis of $L$:
\[
  F_a(B_{\rho}z)
   =\sum_{b=1}^{n}\omega_bF_{a,b}^{(\rho)}(z),
  \qquad
  F_{a,b}^{(\rho)}\in\QQ[z_1,\ldots,z_{10}]_3.
\]
The $\QQ$-span of all component forms $F_{a,b}^{(\rho)}$ has dimension
thirty.  Row reduction in the $220$-dimensional space of cubic monomials
produces a basis
\[
  F_1^{(\rho)},\ldots,F_{30}^{(\rho)}.
\]
Their common zero locus
\[
  \tM_{\rho}
  =V(F_1^{(\rho)},\ldots,F_{30}^{(\rho)})
  \subseteq\PP(V)\cong\PP^9_{\QQ}
\]
is the descended compactified gamma variety, namely the Zariski closure of
the twisted marked moduli space.

Conceptually, the thirty-dimensional relation space is itself stable under
the semilinear descent action and therefore has a thirty-dimensional
$\QQ$-form.  Splitting coefficients after substitution is a concrete way to
compute that form.  An equivalent implementation could compute the fixed
subspace of the relation representation directly.

\section{Searching for a rational point}

Step (vi) asks for
\[
  z=(z_1:\cdots:z_{10})\in\tM_{\rho}(\QQ).
\]
The paper does not prescribe a general rational-point algorithm.  This is the
least canonical and often the most computationally expensive stage.

A basic search enumerates primitive integral vectors of bounded height and
tests the thirty homogeneous cubics.  The authors improve the coordinate
system before searching:

\begin{enumerate}
\item Intersect the fixed space with $\mathcal{O}_L^{10}$, where
$\mathcal{O}_L$ is the maximal order.  This gives a rank-ten $\ZZ$-lattice.
\item Use the Minkowski embedding of $L$ to give the lattice a Euclidean
structure.
\item Apply LLL and use the reduced lattice basis as the projective
coordinate basis of $V$.
\end{enumerate}

This does not change the descended variety; it changes only its coordinates.
In the authors' computations, the reduced coordinates made small-height
rational points much easier to find.

A point on the projective closure need not lie in the open moduli locus of
smooth marked cubic surfaces.  Consequently, every candidate must be passed
through the reconstruction and smoothness tests described below.

\section{Mapping a rational point back to Coble coordinates}

Suppose a rational point $z\in\QQ^{10}\setminus\{0\}$ has been found.  The
maps back to the original coordinates are explicit:
\[
  z
  \xmapsto{\ B_{\rho}\ }
  y=B_{\rho}z\in L^{10}
  \xmapsto{\ J\ }
  x=Jy\in L^{80}.
\]
Choose the forty designated ``positive'' entries from the signed list to
obtain
\[
  (\gamma_1,\ldots,\gamma_{40})\in L^{40}.
\]
The exact fixed-vector construction guarantees that these coordinates obey
the twisted Galois rule.  Compute $P_2,P_4,P_6,P_8,P_{10}$ and then
$A,B,C,D,E$ from the formulas in \cref{sec:clebsch-operational} below.
The resulting values should be fixed by every automorphism of $L/\QQ$; this
is an important exact-arithmetic check.

\section{The pentahedral family and the equation problem}
\label{sec:clebsch-operational}

\subsection{Sylvester pentahedral form}

A general cubic surface over an algebraic closure has a pentahedral
presentation
\[
\begin{cases}
 a_0W_0^3+a_1W_1^3+a_2W_2^3+a_3W_3^3+a_4W_4^3=0,\\
 W_0+W_1+W_2+W_3+W_4=0.
\end{cases}
\]
The five coefficients are unordered and are defined only up to common
scaling.  Let
\[
  \sigma_i=e_i(a_0,\ldots,a_4),
  \qquad 1\leq i\leq5,
\]
be their elementary symmetric functions.

A proper pentahedron corresponds to
\[
  \sigma_5=a_0a_1a_2a_3a_4\neq0.
\]
On this family, the Clebsch invariants are represented by
\[
\begin{aligned}
 A&=\sigma_4^2-4\sigma_3\sigma_5,\\
 B&=\sigma_1\sigma_5^3,\\
 C&=\sigma_4\sigma_5^4,\\
 D&=\sigma_2\sigma_5^6,\\
 E&=\sigma_5^8.
\end{aligned}
\]
These are formulas (3.2) of the paper.

\subsection{Recovering the elementary symmetric functions}

Given a Clebsch vector $[A:B:C:D:E]$ with $E\neq0$, paragraph A.8 chooses a
convenient common scaling of the pentahedral coefficients and gives
\[
 \boxed{
 [\sigma_1,\sigma_2,\sigma_3,\sigma_4,\sigma_5]
 =\left[B,\ D,\ \frac{C^2-AE}{4},\ CE,\ E^2\right].
 }
\]
The weights are consistent: under weighted scaling of
$[A:B:C:D:E]$, the five entries scale with exponents
$2,4,6,8,10$, which is the scaling obtained by multiplying all pentahedral
coefficients by a common factor.

Form the monic quintic
\[
\begin{aligned}
 g(T)
 &=T^5-\sigma_1T^4+\sigma_2T^3-\sigma_3T^2+\sigma_4T-\sigma_5\\
 &=\prod_{i=0}^{4}(T-a_i).
\end{aligned}
\]
The roots are the pentahedral coefficients.  They generally do not lie in
$\QQ$, even though all $\sigma_i$ do.

\subsection{The trace form of Algorithm A.10}

To avoid computing a splitting field, define the degree-five etale algebra
\[
  \cA=\QQ[T]/(g),
  \qquad \theta=T\bmod g.
\]
The algebra may be a field or a product of fields; separability of $g$ is the
essential condition.

Let
\[
  N=\ker\bigl(\Tr_{\cA/\QQ}:\cA\to\QQ\bigr).
\]
Since $\dim_{\QQ}\cA=5$, the trace-zero subspace $N$ has dimension four.
Choose a basis
\[
  c_0,c_1,c_2,c_3\in N.
\]
Set
\[
  \ell(X)=c_0X_0+c_1X_1+c_2X_2+c_3X_3\in\cA[X_0,X_1,X_2,X_3].
\]
Then Algorithm A.10 outputs the cubic form
\[
 \boxed{
  f(X_0,X_1,X_2,X_3)
  =\Tr_{\cA/\QQ}\bigl(\theta\,\ell(X)^3\bigr).
 }
\]
This is the coordinate-free version of the four steps in Algorithm A.10.
Computationally, one reduces $\theta\ell(X)^3$ modulo $g(T)$ and applies the
algebra trace to each polynomial coefficient.

Under the five embeddings of $\cA$ into $\ol\QQ$, $\theta$ maps to
$a_0,\ldots,a_4$ and $\ell$ maps to five conjugate linear forms
$\ell_0,\ldots,\ell_4$.  The trace-zero condition gives
\[
  \ell_0+\cdots+\ell_4=0,
\]
while the trace cubic becomes
\[
  a_0\ell_0^3+\cdots+a_4\ell_4^3.
\]
This explains structurally why the output is the Galois descent of the
pentahedral surface.  No splitting field of $g$ is needed.

To implement the basis computation exactly, write
\[
  t_j=\Tr_{\cA/\QQ}(\theta^j),\qquad 0\leq j\leq4.
\]
A basis of the kernel of the row matrix
\[
  (t_0\ t_1\ t_2\ t_3\ t_4)
\]
gives the coefficient vectors of $c_0,\ldots,c_3$ in the power basis
$1,\theta,\ldots,\theta^4$.

\section{Algorithm 5.1, step by step}

It is useful to separate universal data from data depending on the selected
group and field.

\subsection{Universal precomputation}

Compute once:

\begin{enumerate}
\item a canonical list of the forty gamma invariants and the eighty signed
labels;
\item the permutation representation
$\Pi:W(E_6)\to S_{80}$, or equivalently signed $40\times40$ matrices;
\item ten independent signed gamma sections and the matrix
$J\in\Mat_{80\times10}(\QQ)$ expressing all eighty in that basis;
\item a basis $F_1,\ldots,F_{30}$ of the cubic relation space in the chosen
ten coordinates.
\end{enumerate}

These objects do not depend on $G$ or $L$.

\subsection{Instance-specific computation}

Given $G$, $L$, and $\rho$:

\begin{enumerate}[label=\textbf{Step \arabic*.}]
\item Choose generators $g_1,\ldots,g_s$ of $G$.  Retrieve their permutations
$\Pi(g_i)$ and compute the ten-dimensional matrices $R_{g_i}$.

\item Compute the field automorphisms
$\sigma_i=\rho^{-1}(g_i)$ as matrices in a fixed $\QQ$-basis of $L$.

\item Solve
\[
  R_{g_i}\sigma_i(y)=y
\]
for all generators.  Obtain a $\QQ$-basis of the ten-dimensional fixed space
$V$ and the change-of-basis matrix $B_{\rho}$.

\item Substitute $y=B_{\rho}z$ into the universal thirty cubic relations,
split coefficients in a $\QQ$-basis of $L$, and row-reduce to thirty cubics
over $\QQ$.  These equations define the compactified twist
$\tM_{\rho}\subset\PP^9_{\QQ}$.

\item Search for a rational point $z\in\tM_{\rho}(\QQ)$.

\item Map $z$ back through $B_{\rho}$ and $J$ to obtain the forty gamma
values.  Compute the Clebsch vector $[A:B:C:D:E]$.

\item If $E\neq0$, recover $\sigma_1,\ldots,\sigma_5$, construct $g(T)$,
and apply the trace formula of Algorithm A.10.  Verify the output.  If a
genericity or smoothness check fails, return to the rational-point search.
\end{enumerate}

The mathematical content of the seven steps can be compressed as follows:
\[
\boxed{
\begin{array}{c}
\text{combine the }W(E_6)\text{-representation with the Galois representation}\\
\Downarrow\\
\text{descend Coble's }\PP^9\text{ and its thirty cubic equations}\\
\Downarrow\\
\text{find a rational moduli point}\\
\Downarrow\\
\text{forget the marking and solve the equation problem.}
\end{array}}
\]

\section{Failure loci and post-computation checks}

The algorithm is generic rather than unconditional for every point found on
the compactification.

\subsection{Failure to find a rational point}

A bounded search may fail even when rational points exist.  It is also
possible that the selected twist has no $\QQ$-rational point.  The paper's
remedy is to choose a different Galois extension $L/\QQ$ with the same Galois
group and repeat the descent.

\subsection{The three main degeneracies}

After computing $[A:B:C:D:E]$, test:

\begin{enumerate}
\item $E\neq0$.  Since $E=\sigma_5^8$, this is the proper-pentahedron
condition needed by paragraph A.8.

\item $g(T)$ is separable.  A repeated root means that two pentahedral
coefficients coincide.  This is the Eckardt/nontrivial-automorphism locus,
denoted by $F=0$ in the paper.  On this locus, the Clebsch invariants determine
only the geometric isomorphism class, not the required $\QQ$-form, so the
equation problem is intrinsically ambiguous from these invariants alone.

\item The output cubic is smooth.  Equivalently, its cubic-surface
discriminant $\Delta$ is nonzero.  In practice, compute the singular locus of
the output form or use a discriminant routine.
\end{enumerate}

A rational point on the boundary of $\tM_{\rho}$ will normally be rejected by
one of these checks or by direct verification of the resulting surface.

\subsection{Recommended certification checks}

For an exact implementation, verify all of the following:

\begin{itemize}
\item $\rank(J)=10$ and $P_gJ=JR_g$ for every group generator;
\item the matrices $R_g$ satisfy the group relations;
\item the universal cubic relation space has dimension $30$ and is stable
under $W(E_6)$;
\item the common fixed space has $\QQ$-dimension $10$ and
$\det(B_{\rho})\neq0$ in $L$;
\item the descended cubic space has $\QQ$-dimension $30$;
\item the point found satisfies all descended cubics;
\item the recovered Clebsch quantities are fixed by $\Gal(L/\QQ)$;
\item $g$ is separable and the trace cubic has coefficients in $\QQ$;
\item the final cubic surface is smooth;
\item for a final independent certification, compute the $27$ lines and
compare their Galois permutation group with the desired conjugacy class in
$W(E_6)$.
\end{itemize}

\section{How the universal data can be reconstructed}

Algorithm 5.1 assumes that the signed permutation action, the linear
relations, and the thirty cubic equations are already explicit.  The paper
does not print large tables containing these data.  The accompanying Magma
example reconstructs them computationally.  The following is the underlying
logic.

\subsection{Linear relations by evaluation}

Sample many rational six-point configurations in general position and
evaluate all forty gamma invariants.  The row space of the resulting
point-coordinate matrix has dimension ten.  A basis of that row space gives
a $10\times40$ matrix parametrizing the linear $\PP^9\subset\PP^{39}$; after
adding the negative coordinates, one obtains $J$ for the eighty-coordinate
model.

\subsection{Cubic relations by interpolation}

Express the sampled gamma points in the chosen ten coordinates.  There are
$220$ monomials of degree three in ten variables.  Evaluate all $220$
monomials at sufficiently many sample points.  The kernel of the resulting
evaluation matrix has dimension thirty and supplies the cubic relation
space.

\subsection{The signed $W(E_6)$ action}

Use generators consisting of permutations of the six points and a standard
quadratic Cremona transformation.  Evaluate the gamma vector at a generic
configuration and at its transformed configuration.  After consistent
normalization, match the entries up to sign.  This recovers a signed
permutation matrix on the forty gamma coordinates.  Duplicating signs yields
the ordinary permutation $\Pi(g)\in S_{80}$.

All interpolation and matching computations should be followed by symbolic
or exact evaluation checks on additional configurations.

\section{Notation and data structures at a glance}

\begin{center}
\renewcommand{\arraystretch}{1.2}
\begin{tabular}{>{\raggedright\arraybackslash}p{0.18\textwidth}
                >{\raggedright\arraybackslash}p{0.30\textwidth}
                >{\raggedright\arraybackslash}p{0.42\textwidth}}
\toprule
Symbol & Type & Role \\
\midrule
$G\subseteq W(E_6)$ & finite subgroup & desired permutation group on the $27$ lines \\
$L/\QQ$ & finite Galois extension & realizes the abstract group $G$ \\
$\rho$ & $\Gal(L/\QQ)\xrightarrow{\sim}G$ & pairs field automorphisms with marking changes \\
$\gamma_i$ & forty marked invariants & Coble coordinates on $\tcM$ \\
$\Pi(g)$ & permutation in $S_{80}$ & action on the signed gamma coordinates \\
$J$ & $80\times10$ rational matrix & expresses all signed gammas in a fixed basis \\
$R_g$ & $10\times10$ rational matrix & action on the ten coordinate parameters \\
$V$ & ten-dimensional $\QQ$-space & fixed space defining the descended $\PP^9$ \\
$B_{\rho}$ & $10\times10$ matrix over $L$ & maps rational twist coordinates to original gamma coordinates \\
$F_a^{(\rho)}$ & thirty rational cubics & equations of the compactified twist \\
$[A:B:C:D:E]$ & point of $\PP(1,2,3,4,5)$ & unmarked Clebsch invariant vector \\
$\sigma_i$ & rational numbers & elementary symmetric functions of pentahedral coefficients \\
$g(T)$ & separable quintic & has the pentahedral coefficients as roots \\
$\cA=\QQ[T]/(g)$ & degree-five etale algebra & supports the second Galois descent \\
$\Tr(\theta\ell^3)$ & cubic form over $\QQ$ & final equation of the cubic surface \\
\bottomrule
\end{tabular}
\end{center}

\section{The main conceptual points to retain}

\begin{enumerate}
\item The desired Galois action is encoded by twisting the \emph{marked}
moduli space, not by imposing equations directly on the $27$ lines.

\item Coble's forty invariants depend on the marking.  Their ten-dimensional
span supplies an explicit projective model on which $W(E_6)$ acts linearly.

\item The eighty signed coordinates are a bookkeeping device that turns a
signed permutation action into an ordinary permutation action.

\item The fixed-space equation combines two matrices: a matrix for a field
automorphism of $L$ and a matrix for the corresponding element of $W(E_6)$.
The common fixed space is the rational form of the ambient $\PP^9$.

\item The thirty cubic relations must be descended as a vector space; after a
change of basis over $L$, individual relations need not have rational
coefficients until their coefficient components are row-reduced.

\item A rational point on the twist gives gamma values over $L$, but the
$W(E_6)$-invariant power-sum formulas produce rational Clebsch invariants.

\item The final equation is obtained from a separate degree-five etale algebra
and the trace form $\Tr(\theta\ell^3)$.  This avoids computing the roots of the
pentahedral quintic.

\item The rational-point search is not specified by the theory and is the
principal heuristic component.  LLL reduction of the fixed integral lattice
is the paper's main practical improvement.
\end{enumerate}

\begin{thebibliography}{9}

\bibitem{EJ}
A.-S. Elsenhans and J. Jahnel,
\emph{Moduli spaces and the inverse Galois problem for cubic surfaces},
Trans. Amer. Math. Soc. \textbf{367} (2015), no.~11, 7837--7861;
\href{https://arxiv.org/abs/1209.5591}{arXiv:1209.5591v3}.

\bibitem{EJcode}
J. Jahnel,
\emph{c9\_example.m: a Magma example implementing Algorithm 5.1},
authors' supplementary webpage,
\url{https://wwwuser.gwdguser.de/~jjahnel/Arbeiten/Kub_Fl/c9_example.m}.

\end{thebibliography}

\end{document}
